% Incorporación focal de dependencias anteriores; no sustituye al núcleo común.
% Procedencia detallada: technical/CIERRE_DEPENDENCIAS_ANTERIORES_VI.md.
% Los valores de prueba se publican después de sus operadores; no son entradas.
\section{Coordenadas generadas, acción y normalización electrónica}
\label{vi:dep:seccion}\label{sec:k-registro}

La construcción de las colecciones utiliza resultados de las etapas anteriores:
los caracteres regionales, el registro dodecafásico, la recta de acción y la
sección electrónica. Reunimos aquí las operaciones que los conectan y las
pruebas necesarias para componerlos. Su orden es causal: las regiones y sus
registros se producen sobre APP mediante TRIT y TPK; las evaluaciones reales
se efectúan después; la elección de unidades interviene al realizar las
magnitudes dimensionales.

Las condiciones iniciales y la selección regional se han desarrollado en
la sección anterior a la aritmética de historias. El dominio del libro
incidencial, sus recuentos y la reconstrucción del registro se definen
expresamente a continuación. La selección de una colección concreta de
visitas y trazas es una operación adicional a la evaluación de un libro:
requiere fijar sus multiplicidades, el transporte de frontera y la
cobertura de las trazas admitidas. Las fórmulas de reconstrucción que se
demuestran aquí permiten evaluar y comprobar ese registro; su
invertibilidad no selecciona por sí sola la colección incidencial.

\subsection{Lectores regionales y compatibilidad de la evaluación}
\label{vi:dep:regiones}

La clausura, la propagación y la autoescala son tres regiones de una misma
generación. La región de clausura tiene cinco emisiones y multiplicidad
\(432+4\cdot144=1008\). En el calibre orientado del catálogo, sus emisiones
son
\[
\begin{array}{c|r}
(501,614,498,169,272,272)&432\\
(810,923,870,169,272,272)&144\\
(810,983,810,169,272,272)&144\\
(870,923,810,169,272,272)&144\\
(870,923,810,418,674,521)&144.
\end{array}
\]
La palabra ternaria común es \(010211\); los representantes de propagación
y autoescala son \(201101\) y \(121200\). El censo y la orientación
preceden a las evaluaciones que siguen.

En las cinco posiciones del lector de clausura, la proyección uniforme
\(P_5=\mathbf1\mathbf1^{\mathsf T}/5\), junto con
\(P_5\lambda=\lambda\) y \(\mathbf1^{\mathsf T}\lambda=1\), determina
\(\lambda=\mathbf1/5\). Esta normalización de las posiciones del lector
es distinta de la distribución de semillas indicada en la tabla.
En el régimen elíptico \(J^2=-I\), la multiplicación da
\[
(I+aJ)(I+bJ)=(1-ab)I+(a+b)J.
\]
Siempre que \(1-ab\ne0\), la pendiente compuesta es
\(a\oplus b=(a+b)/(1-ab)\). Dos duplicaciones de \(1/5\) dan
\(5/12\) y \(120/119\). El compensador positivo \(1/q\), \(q>1\),
que devuelve la pendiente a uno satisface
\[
\frac{120/119-1/q}{1+(120/119)/q}=1,
\qquad q=239.
\]
Por tanto, el carácter de clausura de este lector es
\begin{equation}
 \Pi_{\rm cl}=4\left(4\arctan\frac15-\arctan\frac1{239}\right).
 \label{vi:dep:lector-clausura}
\end{equation}
La regla de composición, su compensación y la elección de la rama angular
fijan este valor; la identidad angular convencional lo reconoce como \(\pi\).

El carácter de propagación se realiza en una colección formal conmutativa
\(P_t\), normalizada por \(P_0=1\), \(P'_0=1\) y
\(P_{t+s}=P_tP_s\). Si \(P_t=\sum_{n\ge0}a_nt^n\), la comparación
del coeficiente de \(s^nt\) da
\[
 (n+1)a_{n+1}=a_n,\qquad a_0=1.
\]
Así, \(a_n=1/n!\); la representación a una unidad de parámetro es
\(\Pi_{\rm pr}=\sum_{n\ge0}1/n!\), reconocida como \(e\).
La unidad y el signo del generador pertenecen al lector: la sola ley
semigrupal admitiría otros generadores.

Para la autoescala, TRIT selecciona los modos
\((\Sigma,\Pi,\Pi)\). Sus tres transiciones cíclicas son
\(\Sigma\to\Pi\), \(\Pi\to\Pi\) y \(\Pi\to\Sigma\). Su matriz
de incidencia, en orden areal--volumétrico, es
\begin{equation}
 F=\begin{pmatrix}0&1\\1&1\end{pmatrix}.
 \label{vi:dep:incidencia-hojas}
\end{equation}
El vector propio positivo normalizado como \((1,x)^{\mathsf T}\)
satisface \(x=\lambda\), \(1+x=\lambda x\); por tanto,
\(x^2-x-1=0\). Su raíz positiva es la coordenada de autoescala
\(\varphi\). Esta ecuación reconoce el carácter de la incidencia ya
producida por el calendario.

\begin{proposition}[Evaluación regional mediante intervalos compatibles]
\label{vi:dep:prefijos}
Los tres lectores anteriores disponen de intervalos racionales de anchura
arbitrariamente pequeña. Sus prefijos decimales finitos se determinan por
refinamiento y son compatibles por truncamiento.
\end{proposition}
\begin{proof}
Para \(0<x<1\), dos sumas alternadas consecutivas de
\(\sum_{j\ge0}(-1)^jx^{2j+1}/(2j+1)\) encierran \(\arctan x\).
El primer término omitido acota el error. Esto proporciona intervalos
racionales para \eqref{vi:dep:lector-clausura}.
Si \(E_n=\sum_{j=0}^n1/j!\), entonces, para \(n\ge1\),
\[
 E_n<\Pi_{\rm pr}<E_n+
 \frac{n+2}{(n+1)(n+1)!}.
\]
La cota resulta de mayorar los cocientes sucesivos de la cola por
\(1/(n+2)\). Finalmente, \(F^n\) tiene entradas
\(F_{n-1},F_n,F_{n+1}\), donde \(F_0=0,F_1=1\) y
\(F_{n+1}=F_n+F_{n-1}\). La identidad
\(F_{n+1}^2-F_nF_{n+2}=(-1)^n\) muestra que dos cocientes consecutivos
encierran \(\varphi\) con anchura \(1/(F_nF_{n+1})\).
Se pueden intersectar los intervalos sucesivos para conservar el
encajamiento. Tras el reconocimiento convencional, la irracionalidad de
\(\pi,e,\varphi\) asegura que ninguna coordenada está en una frontera
decimal racional. Para cada profundidad finita, un intervalo
suficientemente estrecho queda en una única celda decimal. Su índice
determina el prefijo y el encajamiento demuestra su compatibilidad.
\end{proof}

\subsection{Coordenatización dodecafásica firmada del estado TPK pleno}
\label{vi:dep:estado-pleno}

La evolución que produce el registro se formula sobre el estado completo del
Topological Phase Kernel. Denotemos por
\(\mathsf{TPKFullState}\) el espacio de estados que conserva conjuntamente
posición APP, régimen TRIT, ruta, hoja, orientación, acarreo, frontera,
memoria, libro incidencial y coordenadas terminales. En una coordenatización canónica,
\[
 X=(X_{\rm APP},X_{\rm TRIT},X_{\rm ruta},X_{\rm mem},
       C_{\rm term},U_{12}^{\rm sgn},X_{\partial},\ldots).
\]
La dinámica de ciento ocho pasos actúa mediante la composición
\begin{equation}
 X_{\rm term}
 =\bigl(\mathcal U_{108}\circ\cdots\circ\mathcal U_1\bigr)(X_{\rm seed}),
 \qquad
 \mathcal U_t=\mathsf{Upd}_t\circ\mathsf{Tra}_t\circ\mathsf{Sel}_t,
 \label{vi:dep:estado-terminal-108}
\end{equation}
donde \(\mathsf{Sel}_t\) aplica el selector trítico de fase,
\(\mathsf{Tra}_t\) transporta hoja, ruta y orientación, y
\(\mathsf{Upd}_t\) incorpora el acarreo y la memoria producidos en el
instante \(t\). Cada factor actúa, por tanto, sobre el estado pleno y la
composición conserva la genealogía de las doce ventanas nonádicas.

El libro incidencial de la sección siguiente es la coordenada
\(\mathcal R_{12}(X_{\rm term})\). Su evaluación por los canales
\(90/120\) define
\[
 \mathcal E_{108}^{90,120}\circ\mathcal R_{12}:
 \mathsf{TPKFullState}\longrightarrow\mathbb Z^{25},
 \qquad
 X\longmapsto C(X)=(b^{90}(X),b^{120}(X),Q(X)).
\]
La proyección dodecafásica firmada es
\[
 \mathcal S_{12}:=\operatorname{pr}_{U_{12}^{\rm sgn}}:
 \mathsf{TPKFullState}\longrightarrow\mathbb Z^{12}.
\]
Sobre la subred compatible producida por la dinámica, las dos coordenatizaciones
satisfacen la identidad operatoria
\begin{equation}
 \mathcal S_{12}
 =\Pi_H\circ\mathcal E_{108}^{90,120}\circ\mathcal R_{12}
 =\Pi_H\circ\operatorname{pr}_{C}.
 \label{vi:dep:identidad-productor-e108}
\end{equation}
Así, la firma dodecafásica se obtiene por proyección y evaluación del mismo
estado que conserva el libro; su dominio incluye la orientación y la memoria
que intervienen en cada recuento. La proyección cruda
\(\rho_{108}:\mathsf{TPKFullState}\to\mathsf{CT108RawProjection}\)
designa, en cambio, la coordenatización que olvida esas coordenadas después de la
producción. Esta separación de tipos fija el orden causal entre generación,
representación firmada y lectura reducida.

\subsection{De las incidencias temporales al registro dodecafásico}
\label{vi:dep:registro}\label{subsec:k-registro-integral}

\begin{definition}[Dominio del libro incidencial orientado]
\label{vi:dep:libro-dominio}
Un libro \(\mathscr L\) contiene una colección finita no vacía de rutas,
sus visitas identificadas y una colección finita de trazas. Cada ruta
conserva una visita por instante durante los primeros \(108\) instantes,
y las visitas posteriores que utilicen sus trazas. Una visita conserva
\[
 v=(\gamma,t,i,j,s,\tau,w,\varepsilon_\partial),
\]
donde \(\gamma\) identifica la ruta, \(t\) su instante, \(1\le i,j\le9\)
su posición APP, \(s\in\{\Sigma,\Pi,0\}\) su hoja,
\(\tau\in\{-1,0,1\}\) su orientación trítica y \(w\) su multiplicidad
entera positiva. La coordenatización incidencial evalúa
\[
 n_v=
 \begin{cases}
 i+j,&s=\Sigma,\\
 ij,&s=\Pi,\\
 9,&s=0.
 \end{cases}
 \qquad
 r_v=1+((n_v-1)\bmod9),\quad q_v=(n_v-r_v)/9.
\]
La lectura \(n_v=9\) fija la convención de umbral de esta coordenatización.
Para \(r_v=9\), la etiqueta \(\varepsilon_\partial\) distingue
corona e interior; no se obtiene únicamente del residuo.

Una traza \(\theta=(v_1,\ldots,v_N)\) recorre instantes consecutivos de
una misma ruta. Conserva una multiplicidad \(u(\theta)\in\mathbb N_{>0}\)
y una unidad de longitud reducida
\(\lambda(\theta)\in\mathbb N_{>0}\), con
\[
 N\lambda(\theta)=120(1+3j(\theta)),\qquad j(\theta)\in\mathbb N_0.
\]
Su canal es
\(\sigma(\theta)=\operatorname{sgn}(\sum_{a=1}^N\tau(v_a))\).
La colección contada en los canales \(120\) comprende las trazas de
signo no nulo. La longitud reducida y el número de instantes quedan
relacionados por \(\lambda\); no se identifican entre sí.
\end{definition}

La definición especifica el dominio de la evaluación, incluidos los
datos que deben transmitirse al lector. Sus multiplicidades y la colección
de trazas pertenecen a la selección dinámica del libro. Una enumeración
del dominio o una comprobación de formato no sustituye esa selección.

La generación conserva un libro \(\mathscr L\) de visitas y trazas.
La coordenatización temporal lo divide en ventanas
\(E_m=\{9m-8,\ldots,9m\}\), \(1\le m\le12\),
equivalentes a las posiciones \(m=\beta+1\) del registro
\(\mathcal R_{12}\). Cada ventana retiene su subruta,
orientación, incremento de frontera e incidencias locales.
Cada visita registra posición APP, hoja aritmética, instante, multiplicidad
y orientación de frontera. En cada evaluación \(n_v\) se conserva la
división \(n_v=r_v+9q_v\), con \(1\le r_v\le9\). En la ventana
\(m\), el lector incidencial expresa
\begin{align}
 A_m&=\sum_{v\in\mathscr L_m}w(v)
      \mathbf1_{\{1,3,5,7\}}(r_v),\notag\\
 C_m&=\sum_{j\ge0}
       \bigl(\nu^-_{120,m}(j)-\nu^+_{120,m}(j)\bigr),\notag\\
 V_m&=\sum_{v\in\mathscr L_m}w(v)
      \mathbf1_{\{9\}}(r_v)\,\varepsilon_{\partial}(v).
 \label{vi:dep:incidencias}
\end{align}
Aquí las trazas de cada ventana tienen soporte finito,
\(\varepsilon_{\partial}=+1\) en la corona y \(-1\) en el interior,
y las etiquetas \(\pm,120\) pertenecen al transporte conservado.
La multiplicidad \(w\) y las trazas \(\nu\) son las obtenidas de ese
transporte; no se reconstruyen escogiendo un valor de \(\alpha\).

En términos de las trazas identificadas, la segunda suma tiene la
expresión finita
\[
 C_m=-\sum_{\theta:\,t(v_1)\in E_m}u(\theta)\sigma(\theta).
\]
En efecto, agrupar por \(j(\theta)\) y por signo produce
\(\nu^-_{120,m}(j)-\nu^+_{120,m}(j)\). Esta igualdad especifica
la operación efectivamente aplicada a cada registro; conserva
separados el canal, la multiplicidad y la longitud.

\begin{proposition}[Actualización y concatenación del libro]
\label{vi:dep:libro-actualizacion}
Fijada la colección de visitas y trazas admitidas por el transporte, sus
recuentos se actualizan mediante incrementos enteros. Al añadir una
visita \(v\) de la ventana \(m\), el incremento de \((A_m,C_m,V_m)\) es
\[
 \left(w(v)\mathbf1_{\{1,3,5,7\}}(r_v),\,0,\,
 w(v)\mathbf1_{\{9\}}(r_v)\varepsilon_\partial(v)\right).
\]
Al completarse una traza \(\theta\) cuyo origen pertenece a \(E_m\),
el incremento es \((0,-u(\theta)\sigma(\theta),0)\).
Para dos segmentos consecutivos, el recuento de su concatenación es
la suma de los recuentos de las visitas y trazas contenidas en cada
segmento, más el de las trazas que cruzan su unión.
\end{proposition}
\begin{proof}
Cada suma de \eqref{vi:dep:incidencias} se indexa por identificadores
de visitas o de trazas, por lo que añadir un identificador añade
exactamente el sumando escrito. Para la concatenación, las visitas
se reparten entre los dos segmentos. Las trazas se particionan en
las contenidas en el primero, las contenidas en el segundo y las
que atraviesan la unión. La aditividad de las sumas enteras prueba
la fórmula. Conservar los identificadores y las prolongaciones
pendientes evita contar dos veces una traza o perder una que aún
no ha terminado. La reducción decimal se aplica después del
recuento; los cocientes de esa reducción permanecen en la memoria.
\end{proof}

La reducción decimal de estos tres recuentos produce
\begin{equation}
 z_m=100[A_m]_{10}+10[C_m]_{10}+[V_m]_{10},
 \qquad 0\le z_m<1000.
 \label{vi:dep:bloque-incidencial}
\end{equation}
Sean \((D_3z)_i=z_i-z_{i+3}\),
\((D_4z)_i=z_i-z_{i+4}\), con índices en \(\ZZ/12\ZZ\), y
\(Q(z)=\sum_i z_i\). El lector conserva
\(\mathcal D(z)=(D_3z,D_4z,Q(z))\). La reducción decimal olvida
múltiplos de diez de cada recuento, mientras el libro original conserva
las visitas y los transportes que los generaron.
La lectura de los eventos de código par sobre esa misma historia,
su descomposición \(12=1+5+6\) y su inversión integral se
desarrollan en la proposición \ref{vi:prop:eventos}. Ese vector
de eventos no sustituye el presente registro de incidencias:
los dominios y las coordenadas conservadas se mantienen distintos.

\begin{proposition}[Conjugación incidencial y reducción decimal]
\label{vi:dep:conjugacion-incidencial}
Sea \(\varrho(m)=13-m\). Consideremos una conjugación del libro
que lleva las visitas y las trazas de la ventana \(m\) a las de
\(\varrho(m)\) mediante biyecciones que conservan sus multiplicidades.
Supongamos que conserva la clase aritmética y la etiqueta de frontera
de las visitas, e intercambia los canales positivo y negativo de las
trazas. Entonces los recuentos de \eqref{vi:dep:incidencias} satisfacen
\begin{equation}
 (A_m^\dagger,C_m^\dagger,V_m^\dagger)
 =(A_{\varrho(m)},-C_{\varrho(m)},V_{\varrho(m)}).
 \label{vi:dep:conjugacion-recuentos}
\end{equation}
Escribiendo \(u_m=[C_m]_{10}\), sus bloques decimales cumplen
\begin{equation}
 z_m^\dagger=z_{\varrho(m)}
       +100\mathbf1_{\{u_{\varrho(m)}\ne0\}}-20u_{\varrho(m)}.
 \label{vi:dep:conjugacion-bloque}
\end{equation}
\end{proposition}
\begin{proof}
Las biyecciones conservan cada sumando de \(A_m\) y \(V_m\).
El intercambio de los canales permuta \(\nu^-\) y \(\nu^+\),
por lo que cambia el signo de \(C_m\). Se obtiene
\eqref{vi:dep:conjugacion-recuentos}. Si \(C=10q+u\),
\(0\le u<10\), la división euclídea de \(-C\) da
\[
 [-C]_{10}=10\mathbf1_{\{u\ne0\}}-u,
 \qquad
 \left\lfloor\frac{-C}{10}\right\rfloor
 =-q-\mathbf1_{\{u\ne0\}}.
\]
En \eqref{vi:dep:bloque-incidencial}, sólo cambia el dígito
intermedio. Su diferencia es
\(10([-C]_{10}-[C]_{10})=100\mathbf1_{\{u\ne0\}}-20u\),
lo que demuestra \eqref{vi:dep:conjugacion-bloque}. La segunda
identidad conserva simultáneamente el cociente entero de la reducción.
\end{proof}

La proposición precisa la acción sobre el libro antes de transportar
su representación decimal. Su aplicación a la reversión de una ruta
conserva también el convenio con que se leen los extremos. Sea
\(\gamma=(x_0,\ldots,x_n)\) una ruta enriquecida y sea \(f\) un
observable entero evaluado antes de cada avance. La ruta inversa
\(\gamma^{-1}=(x_n,\ldots,x_0)\) verifica
\begin{equation}
 \sum_{k=0}^{n-1}f(x_{n-k})
 =\sum_{k=0}^{n-1}f(x_k)+f(x_n)-f(x_0).
 \label{vi:dep:conjugacion-frontera}
\end{equation}
En efecto, el miembro izquierdo suma los índices \(1,\ldots,n\);
el primer sumando del derecho suma \(0,\ldots,n-1\). Su diferencia
es exactamente el término de frontera indicado. La identidad se aplica
a cada observable incidencial con sus datos conservados. Los extremos
se contabilizan en los recuentos enteros antes de reducir módulo diez;
en una ruta cerrada el término se anula, mientras que en una ventana
abierta puede cambiar su bloque. Las trazas que atraviesan una unión
se conservan mediante la partición de la
proposición~\ref{vi:dep:libro-actualizacion}.

Así se conectan los recuentos con la conjugación de rutas de la
sección~\ref{vi:sec:conjugacion}, manteniendo declarada la acción
concreta sobre cada incidencia. La negación de un canal y la reversión
de una ruta son operaciones sobre el libro; su reducción decimal es la
representación que acaba de calcularse.

\begin{proposition}[Imagen integral y reconstrucción del registro]
\label{vi:dep:imagen-integral}
Sean \(b,c\in\ZZ^{12}\), \(q\in\ZZ\), con índices de cero a once.
Definamos
\[
 w_i=c_i-b_{i+1},\qquad P_0=0,\qquad
 P_i=\sum_{j=0}^{i-1}w_j,\qquad T_w=\sum_{i=0}^{11}P_i.
\]
Existe un único \(k\in\ZZ^{12}\) tal que
\(D_3k=b\), \(D_4k=c\), \(Q(k)=q\), si y sólo si
\begin{equation}
 b_i=w_i+w_{i+1}+w_{i+2},\qquad
 \sum_iw_i=0,\qquad q+T_w\equiv0\pmod{12}.
 \label{vi:dep:condiciones-imagen}
\end{equation}
En ese caso,
\begin{equation}
 k_i=\frac{q+T_w}{12}-P_i.
 \label{vi:dep:inversa-incidencial}
\end{equation}
\end{proposition}
\begin{proof}
Si existe \(k\), entonces
\(c_i-b_{i+1}=k_i-k_{i+1}\). La suma cíclica de estas diferencias
es cero, su suma de longitud tres es \(b_i\), y
\(k_i=k_0-P_i\). Al sumar las doce coordenadas,
\(q=12k_0-T_w\), lo que prueba las condiciones necesarias y la fórmula.
Recíprocamente, las condiciones hacen integral la fórmula y la vuelven
cíclica. Sus diferencias de longitud tres son \(b_i\), y las de
longitud cuatro son \(w_i+b_{i+1}=c_i\). Su suma es \(q\).
La diferencia entre dos soluciones tiene todas sus diferencias
consecutivas nulas y suma cero; por ello es nula.
\end{proof}

La reconstrucción es una operación sobre la imagen tipada, no una
selección retrospectiva de sus argumentos. En particular,
\(\mathcal D(z)=\mathcal D(z')\) implica \(z=z'\).
Como la expresión de un entero de \([0,999]\) en tres cifras es única,
dos ternas de recuentos expresan el mismo registro si y sólo si coinciden
componente a componente módulo diez.

El cambio de representación firmado utiliza el bloque
\begin{equation}
 H_4=\begin{pmatrix}
 1&1&1&1\\1&1&-1&-1\\
 1&-1&1&-1\\1&-1&-1&1
 \end{pmatrix},\qquad H_4^2=4I_4.
 \label{vi:dep:hadamard}
\end{equation}
El operador \(H_{12}\) aplica \(H_4\) a las tres órbitas ordenadas
\((1,4,7,10)\), \((2,5,8,11)\), \((3,6,9,12)\).
Sobre la subred
\(\mathcal L_H=\{U\in\ZZ^{12}:H_{12}U\equiv0\pmod4\}\),
las aplicaciones
\begin{equation}
 U=H_{12}K,\qquad K=\tfrac14H_{12}U
 \label{vi:dep:K-reversible}
\end{equation}
son inversas. La congruencia expresa exactamente la integridad de la
segunda aplicación, y \(H_{12}^2=4I\) prueba ambas composiciones.

El registro firmado también se calcula directamente de las coordenadas
transversales, antes de reconstruir \(K\). Con índices de cero a once,
sean, para \(r=0,1,2\),
\[
 S_r=\sum_{j=0}^3c_{r+3j},\quad
 a_0=\frac{q+2S_0+S_1}{3},\quad
 a_1=a_0-S_0,\quad a_2=a_1-S_1,
 \qquad d_{r,j}=b_{r+3j}.
\]
La aplicación armónica expresa por órbitas
\begin{equation}
 \Pi_H^{(r)}(b,c,q)=
 (a_r,d_{r,1}-d_{r,3},d_{r,0}+d_{r,2},d_{r,0}-d_{r,2}).
 \label{vi:dep:PiH}
\end{equation}
Para probar la fórmula sobre la imagen integral, sea \(k\) su
preimagen. Si \(s_r=\sum_jk_{r+3j}\), entonces
\(S_r=s_r-s_{r+1}\) y \(q=s_0+s_1+s_2\), con índices de órbita
módulo tres. Resolver da \(s_r=a_r\). Para las cuatro coordenadas
\((x_0,x_1,x_2,x_3)\) de una órbita,
\(d_j=x_j-x_{j+1}\). Sus tres combinaciones en
\eqref{vi:dep:PiH} son respectivamente
\(x_0+x_1-x_2-x_3\), \(x_0-x_1+x_2-x_3\) y
\(x_0-x_1-x_2+x_3\). Junto con \(a_r=\sum_jx_j\), son
exactamente \(H_4(x_0,x_1,x_2,x_3)^{\mathsf T}\).
La composición efectiva es así
\(\mathscr L\to(b,c,q)\to\Pi_H(b,c,q)=U\to H_{12}U/4=K\).
La preimagen utilizada en esta demostración verifica la identidad;
no es un argumento de la aplicación \(\Pi_H\).

Para el estado terminal de \eqref{vi:dep:estado-terminal-108}, la coordenada
producida antes de la representación firmada es
\begin{align}
 b^{90}={}&(-495,-116,-684,108,601,-90,
              -173,-88,313,560,-397,461),\notag\\
 b^{120}={}&(-425,-281,-481,671,-255,30,
               475,-543,680,251,6,-128),\notag\\
 Q_{\rm TPK}={}&6263.
 \label{vi:dep:canales-terminales}
\end{align}
Al sustituir estos veinticinco campos en \eqref{vi:dep:PiH}, las tres
órbitas firmadas son
\begin{align*}
 U_{\rm term}^{(1)}&=(2378,-452,-668,-322),\\
 U_{\rm term}^{(2)}&=(1406,998,-204,-28),\\
 U_{\rm term}^{(3)}&=(2479,-551,-371,-997).
\end{align*}
La intercalación por columnas produce
\begin{equation}
 U_{\rm term}=\mathcal S_{12}(X_{\rm term})
 =(2378,1406,2479,-452,998,-551,-668,-204,-371,-322,-28,-997).
 \label{vi:dep:U-terminal}
\end{equation}
La transformada orbital \(H_{12}/4\) está definida sobre esta subred por
la divisibilidad exacta de cada bloque de cuatro. En consecuencia, la cadena
\begin{equation}
 X_{\rm term}\xrightarrow{\mathcal R_{12}}\mathscr L(X_{\rm term})
 \xrightarrow{\mathcal E_{108}^{90,120}}C_{\rm term}
 \xrightarrow{\Pi_H}U_{\rm term}
 \xrightarrow{H_{12}/4}K
 \label{vi:dep:cadena-terminal-completa}
\end{equation}
es una composición tipada desde el estado pleno hasta su sello periódico.
Las proposiciones~\ref{vi:dep:libro-actualizacion} y
\ref{vi:dep:imagen-integral} proporcionan, respectivamente, la actualización
finita del libro y la inversa integral de la última representación.

El registro generado proporciona el testigo ordenado
\begin{equation}
 K=(234,543,140,729,659,824,621,058,914,794,146,601).
 \label{vi:dep:K-testigo}
\end{equation}
El cero inicial de \(058\) forma parte de la lectura por bloques.
Para comprobarlo mediante la proposición~\ref{vi:dep:imagen-integral}, sus diferencias
y suma son
\[
\begin{split}
b={}&(-495,-116,-684,108,601,-90,-173,-88,313,560,-397,461),\\
c={}&(-425,-281,-481,671,-255,30,475,-543,680,251,6,-128),\\
q={}&6263.
\end{split}
\]
En las tres órbitas anteriores, el registro firmado es
\[
 \begin{aligned}
 U^{(1)}&=(2378,-452,-668,-322),\\
 U^{(2)}&=(1406,998,-204,-28),\\
 U^{(3)}&=(2479,-551,-371,-997).
 \end{aligned}
\]
Aplicar \(H_4/4\) a estos tres vectores reproduce
\eqref{vi:dep:K-testigo}. Estas identidades verifican la representación de
las incidencias; su génesis sigue siendo la del libro
\eqref{vi:dep:incidencias}.

\subsection{\texorpdfstring{Representación periódica de \(K\) y normalización de \(\alpha\)}{Representación periódica de K y normalización de alfa}}
\label{vi:dep:alpha}

Fijemos \(B=1000\), el origen de fase y las doce posiciones. Escribamos
\[
 N_K=\sum_{j=1}^{12}K_jB^{12-j},\qquad
 \kappa=\frac{N_K}{B^{12}-1}.
\]
Ésta es la evaluación periódica exacta del registro. La multiplicación
por \(B^{12}-1\) y la división euclídea sucesiva recuperan sus doce
bloques, incluido cualquier cero inicial. Para
\(\rho K=(K_2,\ldots,K_{12},K_1)\), la misma concatenación da
\(\kappa(\rho K)=B\kappa(K)-K_1\). Por tanto, la reversibilidad
del registro conserva un origen de fase especificado.

\begin{proposition}[Período mínimo de la representación racional]
\label{vi:dep:periodo-minimo-k}
La prolongación periódica del registro \eqref{vi:dep:K-testigo}
tiene período mínimo de \(12\) bloques en base \(1000\) y de
\(36\) cifras en base \(10\).
\end{proposition}
\begin{proof}
Los doce componentes de \(K\) son distintos. Cualquier traslación
cíclica no nula menor que doce identificaría dos componentes si fuera
un período. Por tanto, el período mínimo de bloques es doce.

La concatenación de sus bloques de tres cifras, incluido \(058\),
produce un período decimal de longitud \(36\). Esta expansión es
canónica: no es eventualmente constante en la cifra \(9\).
Sea \(d\) su período decimal mínimo. Entonces \(d\mid36\),
porque los desplazamientos que preservan una palabra periódica forman
un subgrupo cíclico. Al agrupar las cifras de tres en tres, el período
mínimo de bloques es el menor \(k>0\) tal que \(d\mid3k\), es
decir, \(d/\gcd(d,3)\). Ningún divisor propio de \(36\)
satisface \(d/\gcd(d,3)=12\). Se sigue que \(d=36\).
\end{proof}

Este período pertenece a la coordenada racional del registro, con el
origen de fase fijado. La composición siguiente conserva además las
coordenadas regionales y los cocientes de normalización de su salida.

Sean \(P_j,E_j,\Phi_j\) los bloques fraccionarios de los tres
caracteres regionales ya construidos, y extendamos \(K_j\) con período
doce. La composición de registros es
\begin{equation}
 Z_j=P_j+E_j-\Phi_j-K_j,
 \qquad
 A_j=Z_j+c_{j+1}-Bc_j.
 \label{vi:dep:normalizacion-alpha}
\end{equation}
El cociente entero \(c_j\) transporta la continuación hacia la posición
observada. En una ventana finita se fija primero \(c_{N+1}\), y luego
\[
 c_j=\left\lfloor\frac{Z_j+c_{j+1}}B\right\rfloor,
 \qquad A_j=Z_j+c_{j+1}-Bc_j.
\]
La división euclídea determina unívocamente \(0\le A_j<B\).
Como \(-1998\le Z_j\le1998\), el conjunto
\(\{-2,-1,0,1,2\}\) es estable para estos cocientes.

\begin{theorem}[Composición exacta a profundidad ilimitada]
\label{vi:dep:alpha-completa}
La evaluación completa de la normalización, en la convención decimal
canónica, es
\begin{equation}
 \alpha_A=\pi+e-\varphi-4-\kappa.
 \label{vi:dep:alpha-A}
\end{equation}
Existe una única expansión canónica \((A_j)\), con bloques entre cero
y \(B-1\), y una sucesión acotada de cocientes enteros que satisfacen
\eqref{vi:dep:normalizacion-alpha} y \(c_1=0\).
\end{theorem}
\begin{proof}
Las partes enteras regionales son \(3,2,1\). Por convergencia absoluta,
\[
 y:=\sum_{j\ge1}Z_jB^{-j}
   =(\pi-3)+(e-2)-(\varphi-1)-\kappa.
\]
Los primeros bloques son \(141,718,618,234\), de modo que
\(Z_1=7\). La suma de la cola tiene valor absoluto a lo sumo
\(1998\sum_{j\ge2}B^{-j}=0.002\); en particular \(0<y<1\).
Tomemos la expansión canónica \(y=\sum A_jB^{-j}\), eligiendo la
expansión terminal y no la eventualmente constante \(B-1\) si se
presenta una frontera. Pongamos
\[
 c_j=\sum_{r\ge0}(Z_{j+r}-A_{j+r})B^{-r-1}.
\]
La igualdad de los valores completos muestra que cada \(c_j\) es el
negativo de una suma finita de enteros ponderados por potencias de
\(B\); es por tanto entero, y \(c_1=0\). Separar el primer término
da \(Bc_j=Z_j-A_j+c_{j+1}\). Las cotas de las colas hacen acotada
la sucesión. Más precisamente, la cola de \(Z\) está en \([-2,2]\)
y la cola canónica de \(A\) en \([0,1)\), por lo que
\(-3<c_j\le2\). Al ser entero, \(c_j\) pertenece al conjunto
\(\{-2,-1,0,1,2\}\) indicado.
Recíprocamente, la suma de la recurrencia hasta \(N\) da
\[
 \sum_{j=1}^N A_jB^{-j}
 =\sum_{j=1}^N Z_jB^{-j}-c_1+c_{N+1}B^{-N}.
\]
La acotación anula el último término en el límite. La unicidad de la
expansión canónica de \(y\) determina \(A\), y la fórmula de colas
determina entonces todos los \(c_j\).
\end{proof}

La evaluación finita \(\kappa_{36}=N_K/B^{12}\) difiere de
\(\kappa\) en \(B^{-12}\kappa\). Además, imponer
\(c_{13}=0\) en una ventana de doce bloques no sustituye el cociente
de la continuación completa. En el registro expresado, la continuación
da \(c_{13}=-1\): el duodécimo bloque es \(662\), mientras la
ventana cortada con cociente terminal cero da \(663\). La evaluación
completa comienza por
\[
 \alpha_A=0.007297352569283800997285105472380662999\ldots.
\]
El período de \(K\) no impone período a \((A_j)\), pues en
\eqref{vi:dep:normalizacion-alpha} también intervienen los tres
registros regionales y los cocientes enteros.

En las fórmulas siguientes, \(\alpha\) designa esta salida posicional
completa \(\alpha_A\). La vía analítica de vacancias tiene su operador
propio; un polinomio truncado de orden nueve y la raíz de ese polinomio
no se identifican aquí con el operador completo. Las composiciones de
masas que siguen no requieren promover aquella aproximación a una
igualdad entre las dos vías infinitas.

\subsection{Construcción tridimensional del prefactor electrónico}
\label{vi:dep:prefactor}

Sea \(\mathcal A_3=\{1,\ldots,9\}^3\), con medida uniforme, y
\(\rho_9(n)=1+((n-1)\bmod9)\). En este dominio actúan
\[
 \Sigma(x)=\rho_9(x_1+x_2+x_3),\qquad
 \Pi(x)=\rho_9(x_1x_2x_3).
\]
En \(\ell^2(\mathcal A_3)\), sean \(P\) y \(Q\) las medias
condicionales sobre las fibras de \(\Sigma\) y \(\Pi\),
respectivamente. Cada una es idempotente y autoadjunta: dentro de cada
fibra sustituye un vector por su componente constante, y conserva
ortogonalmente esa componente. Son, pues, proyectores ortogonales.

Las fibras aditivas tienen cardinal \(81\). Las multiplicativas tienen
cardinales
\[
m=(36,36,108,36,36,108,36,36,297).
\]
La incidencia conjunta
\(N_{ab}=\#\{x:\Sigma(x)=a,\Pi(x)=b\}\) es \(9N_0\), donde
\begin{equation}
 N_0=\begin{pmatrix}
0&1&1&0&1&1&0&1&4\\1&0&1&1&0&1&1&0&4\\
1&0&2&0&1&2&0&0&3\\0&1&1&0&1&1&0&1&4\\
1&0&1&1&0&1&1&0&4\\0&0&2&1&0&2&0&1&3\\
0&1&1&0&1&1&0&1&4\\1&0&1&1&0&1&1&0&4\\
0&1&2&0&0&2&1&0&3
\end{pmatrix}.
 \label{vi:dep:incidencia-cubica}
\end{equation}
Esta matriz se obtiene contando las \(729\) ternas. También se
reproduce, sin redondeo, con
\(N_{ab}=\sum_{i,j,k=1}^9
\mathbf1_{\{a\}}(\rho_9(i+j+k))
\mathbf1_{\{b\}}(\rho_9(ijk))\).
En las bases normalizadas de funciones constantes por fibra, el
entrelazamiento tiene matriz \(C_{ab}=N_{ab}/\sqrt{81m_b}\).

\begin{proposition}[Norma de la no conmutatividad de las dos hojas]
\label{vi:dep:norma-PQ}
Los proyectores anteriores satisfacen
\begin{equation}
 \|[P,Q]\|=\frac{\sqrt3}{4}.
 \label{vi:dep:prefactor-valor}
\end{equation}
\end{proposition}
\begin{proof}
En el espacio de las nueve fibras aditivas, \(PQP\) tiene matriz
\(M=CC^{\mathsf T}\). La multiplicación de las matrices de
\eqref{vi:dep:incidencia-cubica} da autovalor uno sobre el vector
constante y, sobre los dos generadores no constantes,
\[
 \frac14\quad\hbox{sobre}\quad(1,-1,0,1,-1,0,1,-1,0),
 \qquad
 \frac7{132}\quad\hbox{sobre}\quad(1,1,-2,1,1,-2,1,1,-2).
\]
Sobre las coordenadas \(\{3,6,9\}\) con suma cero actúa por
\(1/18\), con multiplicidad dos. Sobre cada uno de los conjuntos
\(\{1,4,7\}\) y \(\{2,5,8\}\), los vectores de suma cero
pertenecen al núcleo; dan cuatro dimensiones. Estos subespacios
ortogonales suman las nueve dimensiones.

Para un autovector unitario \(u\in\operatorname{ran}P\) de
\(PQP\) con autovalor \(0<\lambda<1\), el vector
\(v=(Qu-\lambda u)/\sqrt{\lambda(1-\lambda)}\) es unitario y
ortogonal a \(\operatorname{ran}P\). En la base \((u,v)\),
\[
 P=\begin{pmatrix}1&0\\0&0\end{pmatrix},\quad
 Q=\begin{pmatrix}\lambda&\sqrt{\lambda(1-\lambda)}\\
 \sqrt{\lambda(1-\lambda)}&1-\lambda\end{pmatrix}.
\]
El conmutador tiene norma \(\sqrt{\lambda(1-\lambda)}\).
Los bloques correspondientes a \(\lambda=0,1\) conmutan. De los
autovalores calculados, el máximo se alcanza en \(\lambda=1/4\),
pues \(\lambda(1-\lambda)\) es creciente en \([0,1/2]\).
Su raíz cuadrada es \(\sqrt3/4\).
\end{proof}

En el bloque extremal, \(S=2P-I\) y
\(J=4[P,Q]/\sqrt3\) satisfacen
\[
 S^2=I,\qquad J^2=-I,\qquad SJ=-JS.
\]
Las cuatro matrices \(I,S,J,SJ\) son linealmente independientes,
de modo que realizan \(\mathrm{Cl}_{1,1}\simeq M_2(\RR)\).
Los elementos \(n_\pm=(S\pm J)/2\) cumplen
\(n_\pm^2=0\) y \(n_+n_-+n_-n_+=I\), por expansión directa.
El prefactor de la ecuación electrónica procede así de dos evaluaciones
de APP en dimensión tres y de su conmutador, antes de cualquier masa
de referencia.

\subsection{Lectura de frontera y composición electrónica basal}
\label{vi:dep:basal}

La incidencia \eqref{vi:dep:incidencia-hojas} tiene vector propio
positivo \(r=(1,\varphi)^{\mathsf T}\). La normalización
\(p_{ij}=F_{ij}r_j/(\varphi r_i)\) produce
\[
 P_{\rm av}=\begin{pmatrix}0&1\\\varphi^{-2}&\varphi^{-1}\end{pmatrix},
 \qquad \mu=\frac{(1,\varphi^2)}{1+\varphi^2}.
\]
La identidad \(\varphi^{-2}+\varphi^{-1}=1\) prueba que las filas
suman uno. Como \(F\) es simétrica, los pesos proporcionales a
\(r_i^2\) satisfacen balance detallado; por ello \(\mu P_{\rm av}=\mu\).
La incidencia es irreducible y tiene una autorretención, de modo que
esta distribución estacionaria es única. La frecuencia cronológica
\((1/3,2/3)\) cuenta los tres instantes del calendario; \(\mu\)
pondera sus trayectorias admisibles de memoria uno.

El operador que conserva las hojas, fija la normalización volumétrica
en uno y marca la areal por el carácter de clausura es
\(D_\pi=\operatorname{diag}(\pi,1)\). Por inducción,
\[
 F^n=\begin{pmatrix}F_{n-1}&F_n\\F_n&F_{n+1}\end{pmatrix}.
\]
El cociente de retornos diagonales marcados es
\[
 \Theta_n=\pi\frac{F_{n-1}}{F_{n+1}}
 \longrightarrow\frac{\pi}{\varphi^2}.
\]
La convergencia se deduce del autovalor dominante \(\varphi\), o de
los cocientes consecutivos ya acotados. El lector
\(\mathcal R(x,y)=x/y^2\), bajo el intercambio de argumentos,
expresa también
\begin{equation}
 \rho_A=\frac{\pi}{\varphi^2},\qquad
 \rho_E=\frac{\varphi}{\pi^2},\qquad
 \varphi^3\rho_A^2\rho_E=1.
 \label{vi:dep:dos-orientaciones}
\end{equation}
La última igualdad se comprueba multiplicando las fracciones.
Relaciona dos orientaciones del mismo lector y no selecciona los
restantes coeficientes de la ecuación electrónica.

El retorno central tipado tiene longitud \(27\), con soporte
\(\ZZ/9\ZZ\times\FF_3\). Las cinco posiciones de la microfibra
se insertan, en el calibre regional fijado, como
\((4,0),(5,0),(6,0),(6,1),(6,2)\).
El complemento tiene \(27-5=22\) posiciones. Su incidencia con los
tres estados de fase trítica tiene \(5\cdot3=15\) pares.
Así se obtienen los enteros \(-22\) y \(15\) utilizados por el
lector electrónico; la orientación fija el signo del primer término.
La composición cúbica utiliza tres copias de la coordenada de cierre,
por lo que su evaluación es \(\alpha^3\). El soporte, la inserción
regional y esta regla de composición son datos operativos distintos;
un recuento aislado no sustituye su declaración.

El defecto global y su retorno nonádico se expresan mediante
\begin{equation}
 d_4=\frac{\pi-e\log\pi}{270},\qquad
 \Delta_4=d_4\left(1+\frac{\pi}{729}\right).
 \label{vi:dep:delta4}
\end{equation}
El denominador \(270\) pertenece a la corona sin su unidad de
clausura; \(729=9^3\) al soporte cúbico. La función
\(f(x)=x-e\log x\) alcanza su mínimo cero en \(x=e\), pues
\(f'(x)=1-e/x\) cambia allí de signo. Como \(\pi\ne e\),
se obtiene \(d_4>0\) y \(\Delta_4>0\).

\begin{definition}[Coordenada electrónica basal]
\label{vi:dep:electron-basal}
La composición de los lectores anteriores es
\begin{equation}
 \varepsilon_e^{(0)}=
 \frac{\sqrt3}{4}
 \left[\exp\left(\frac{\varphi}{\pi^2}\right)-22\alpha^3\right]
 (1+15\Delta_4).
 \label{vi:dep:formula-basal}
\end{equation}
\end{definition}
La exponencial actúa sobre la orientación \(\rho_E\) de
\eqref{vi:dep:dos-orientaciones}; la lectura de vacancia corrige esa
componente y el último factor conserva el retorno del defecto global.
Para \(0<\alpha<0.01\), todos los factores son positivos:
\(\exp(\varphi/\pi^2)>1\), \(22\alpha^3<22\cdot10^{-6}<1\)
y \(\Delta_4>0\). Ésta es una coordenada adimensional. Su presencia
en una coordenatización de masa se define después de la construcción de la escala.

\subsection{Escala determinantal y constante reducida de Planck}
\label{vi:dep:planck}

El bloque local \eqref{vi:dep:hadamard} tiene determinante \(-16\).
Su forma normal de Smith es \(\operatorname{diag}(1,2,2,4)\).
En efecto, el máximo divisor de las entradas es uno; los menores de
orden dos son pares y uno vale dos en valor absoluto. Los menores de
orden tres son, por la identidad de la adjunta, de valor absoluto cuatro,
y el determinante tiene valor absoluto dieciséis. Los divisores
determinantales son \(1,2,4,16\), cuyas razones sucesivas dan los
cuatro factores de Smith. Por consiguiente,
\begin{equation}
 G_H=\operatorname{coker}H_4
 \simeq\ZZ/2\ZZ\oplus\ZZ/2\ZZ\oplus\ZZ/4\ZZ,
 \qquad |G_H|=16.
 \label{vi:dep:conucleo}
\end{equation}

El lector normado local multiplica la misma coordenada \(\alpha\)
una vez por cada clase de ese conúcleo, y aplica una vez la coordenada
de autoescala. Su definición produce
\begin{equation}
 \Sigma_H=\varphi\prod_{g\in G_H}\alpha
         =\varphi\alpha^{16}.
 \label{vi:dep:sigma-H}
\end{equation}
El orden del conúcleo determina el exponente de este lector. La forma
normal de Smith y la elección del lector normado cumplen funciones
distintas: la primera no convierte cualquier funcional del bloque en
la expresión \eqref{vi:dep:sigma-H}.

\begin{proposition}[Década de acción]
\label{vi:dep:decada}
La salida \eqref{vi:dep:sigma-H} satisface
\(10^{-34}<\Sigma_H<10^{-33}\); su valoración decimal es \(-34\).
\end{proposition}
\begin{proof}
Las evaluaciones regionales y posicionales dan los intervalos
\(1.618<\varphi<1.619\) y \(0.007297<\alpha<0.007298\).
Por positividad,
\[
 \frac{1618\,7297^{16}}{10^{99}}
 <\Sigma_H<
 \frac{1619\,7298^{16}}{10^{99}}.
\]
Las desigualdades enteras
\(1618\,7297^{16}>10^{65}\) y
\(1619\,7298^{16}<10^{66}\) prueban el resultado sin utilizar una
constante de acción externa.
\end{proof}

La década y la mantisa son representaciones de etapas diferentes.
Antes de construir la segunda, el lector regional selecciona el retorno
persistente de la microfibra de clausura. La proyección
\(p_{\rm ret}(u_1,\ldots,u_6)=(u_4,u_5,u_6)\) aplicada a las
cinco emisiones de \ref{vi:dep:regiones} produce el bloque
\(b_\pi=(169,272,272)\) cuatro veces y
\(b_\pi^-=(418,674,521)\) una vez. Estas multiplicidades cuentan
filas regionales, no las multiplicidades de semillas \(432,144\).

\begin{proposition}[Selector equivariante del retorno regional]
\label{vi:dep:selector169}
Entre las reglas que conservan las permutaciones de las cinco regiones,
seleccionan el bloque terminal de multiplicidad máxima, son equivariantes
bajo permutaciones de sus tres posiciones y retienen una posición fija
por el estabilizador del bloque, la coordenada seleccionada es única
y vale \(\rho_\pi=169\).
\end{proposition}
\begin{proof}
La proyección de las cinco filas muestra que el bloque de multiplicidad
máxima es único y es \(b_\pi\). Su estabilizador en
\(\mathfrak S_3\) intercambia las dos posiciones de valor \(272\)
y fija únicamente la posición de valor \(169\). Retener esa posición
es compatible con cualquier permutación simultánea de coordenadas.
Equivalentemente, el valor se expresa sin elegir una posición inicial
como
\(\operatorname{sing}(b_\pi)=\sum_j(b_\pi)_j
-2\operatorname{moda}(b_\pi)=169\).
\end{proof}

La graduación por longitud expresa una palabra de longitud \(n\)
mediante el carácter \(\chi_z([w])=z^{|w|}\). La concatenación
suma longitudes, por lo que
\(\chi_z([uv])=\chi_z([u])\chi_z([v])\). Como las emisiones
regionales pertenecen a \(\mathcal U_6\), el impulso del retorno
aparece como \(\rho_\pi z^6\). El grado seis corresponde a la
longitud de la emisión, no a la longitud de un paseo cerrado que
recorra varios anclajes del grafo.

El peso de prolongación se obtiene en la línea determinante del
transporte bicapa. Para una involución real \(R\) en dimensión dos,
con \(\operatorname{tr}R=0\), sean
\[
 x_A=\frac{50\pi}{9}\alpha,\qquad
 \mathcal T(x_A,y)=\exp(-x_AI_2-yR).
\]
La coordenada orientada \(y\) todavía es formal. En una base propia
de \(R\), las dos entradas son \(e^{-x_A-y}\) y
\(e^{-x_A+y}\). Por tanto,
\[
 \bigwedge^2\mathcal T=\det\mathcal T
 =e^{-2x_A}=D_A.
\]
La dependencia en \(y\) se cancela. La multiplicatividad por
concatenación determina el carácter de prolongación \(k\mapsto D_A^k\):
cualquier carácter con valor elemental \(D_A\) debe tener ese valor
en \(k=1+\cdots+1\). Esta construcción precede a la evaluación
de la coordenada angular conjugada.

El lector regional de orden cinco emplea los autovalores \(9,5\)
del bloque central \(\bigl(\begin{smallmatrix}7&2\\2&7\end{smallmatrix}\bigr)\),
la órbita de cuatro posiciones y el calendario de doce ventanas.
Su regla graduada fija
\begin{align}
 H_5={}&\frac12\sqrt{\frac{1000\alpha}{\varphi}}-\alpha
       +\frac9{16}\alpha^2-\frac59\alpha^3
       +\frac7{48}\alpha^4-\frac1{54}\alpha^5,\notag\\
 D_A={}&\exp\left(-\frac{100\pi\alpha}{9}\right),\notag\\
 \mathcal C_\pi={}&169\alpha^6+
                  \frac{D_A\alpha^7}{1-D_A\alpha},\notag\\
 \eta_{\rm ret}={}&H_5-\frac{90}{\pi}\mathcal C_\pi.
 \label{vi:dep:mantisas}
\end{align}
En esa regla, los cocientes son
\(9/16=9/4^2\), \(5/9\), \(7/48=(\operatorname{tr}B_c/2)/(4\cdot12)\)
y \(1/54=1944/104976\). Su asignación a los grados y los signos
forma parte del lector graduado. El entero \(169\) es la representación
del selector demostrado en la proposición \ref{vi:dep:selector169}.
Su dato regional es anterior a la evaluación de \(\alpha\); el
lector analítico lo compone después con la coordenada de cierre.

La cola de \(\mathcal C_\pi\) tiene una construcción por renovación.
El impulso \(\rho_\pi z^6\) y la continuación estricta pertenecen
a sumandos distintos. Después de registrar una sola vez el impulso,
la continuación se normaliza por la unidad de la línea determinante;
cada nuevo paso la multiplica por \(D_A\). La orientación del
calibre regional coloca este retorno en el sector compensador que
se sustrae del carácter de quinto grado.
Con parámetro formal \(z\), la regla
\begin{equation}
 R(z)=D_Az^7+D_AzR(z)
 \label{vi:dep:renovacion}
\end{equation}
determina de modo único
\(R(z)=D_Az^7/(1-D_Az)\). Se puede probar por comparación de
coeficientes o por inversión de la serie formal \(1-D_Az\).
La normalización del primer retorno, además de la razón recurrente,
es esencial para esta unicidad. Como \(0<D_A<1\) y \(\alpha<1\),
la evaluación en \(z=\alpha\) converge absolutamente.

\begin{proposition}[Positividad y razón de retorno de acción]
\label{vi:dep:positividad-accion}
En la rama generada, \(0<\eta_{\rm ret}<H_5\). La razón
\begin{equation}
 R_{\rm act}=\frac{H_5}{\eta_{\rm ret}}
 \label{vi:dep:Ract}
\end{equation}
es positiva, mayor que uno y adimensional.
\end{proposition}
\begin{proof}
Los límites \(0.007<\alpha<0.008\) y \(\varphi<1.619\) dan
\(\tfrac12\sqrt{1000\alpha/\varphi}>1\). La suma del valor absoluto
de los términos negativos restantes de \(H_5\) es menor que
\(0.008001\), por lo que \(H_5>0.99\).
Como \(D_A<1\), \(\pi>3\) y \(\alpha<0.008\),
\[
 0<\frac{90}{\pi}\mathcal C_\pi
 <30\left(169(0.008)^6+
                  \frac{(0.008)^7}{1-0.008}\right)<10^{-6}.
\]
Se deduce \(0<\eta_{\rm ret}<H_5\). La razón de dos coordenadas
de acción en la misma base es adimensional y no cambia al reescalar
esa base.
\end{proof}

Sea \(\mathcal U_S\) una base de la recta de acción. Las dos
secciones, anterior y posterior al retorno, son
\begin{equation}
 \hbar_{\rm pre}=H_5\,10^{-34}\mathcal U_S,
 \qquad
 \hbar_{\rm ret}=\eta_{\rm ret}\,10^{-34}\mathcal U_S.
 \label{vi:dep:hbar}
\end{equation}
Su lectura física es la construcción propuesta de la constante reducida
de Planck, con la sección de retorno explicitada. La unidad
\(\mathcal U_S\mapsto\mathrm{J\,s}\) es una realización metrológica
posterior. El coeficiente de \(\Sigma_H\) dentro de su década no es
\(H_5\): el determinante expresa la escala y el lector regional
expresa la coordenada que ocupa esa escala. Para cada sección,
\(h=2\pi\hbar\) expresa la acción de una vuelta completa.

\subsection{Coordenadas angulares y reloj de la realización absoluta}
\label{vi:dep:reloj}

El par angular utiliza las coordenadas adimensionales ya generadas:
\begin{equation}
 A_{\deg}=1000\alpha,\qquad
 C^*_{\deg}=2(\eta_{\rm ret}+\alpha),\qquad
 x=\frac{\pi}{180}A_{\deg},\quad
 y=\frac{\pi}{180}C^*_{\deg}.
 \label{vi:dep:angulos}
\end{equation}
La primera operación desplaza un bloque la expansión en base mil.
La segunda es la representación bidireccional de la coordenada regional:
equivalentemente,
\(y=\pi(H_5+\alpha)/90-\mathcal C_\pi\).
Esta igualdad se deduce sustituyendo \eqref{vi:dep:mantisas}; no
introduce una regla adicional de selección del contraángulo.
La coordenatización de una vuelta tiene \(360\) unidades angulares o \(2\pi\)
radianes. Así, para una coordenada levantada,
\[
 \hbar\theta_{\rm rad}
 =h\frac{\theta_{\deg}}{360}.
\]
El factor \(10^{-34}\) permanece en la recta de acción y no se
inserta en las coordenadas angulares.

Los canales positivos son
\[
 q_+=e^{-x-y},\qquad q_-=e^{-x+y}.
\]
Se recupera la pareja mediante
\(x=-\tfrac12\log(q_+q_-)\),
\(y=\tfrac12\log(q_-/q_+)\). El intercambio de canales conserva
\(x\) y cambia el signo de \(y\). El operador con proyectores
orientados \(P_\pm\),
\(T=q_+P_++q_-P_-\), conserva esas etiquetas al expresar sus dos
autovalores.

Para una unidad temporal interna \(t_0>0\), el cierre de \(108\)
pasos determina
\begin{equation}
 \omega_0=\frac{2\pi}{108t_0},\qquad
 E_{\rm clk}=\hbar_{\rm ret}\omega_0,
 \qquad m_{\rm clk}=\frac{E_{\rm clk}}{c_{\rm int}^{2}}.
 \label{vi:dep:reloj-absoluto}
\end{equation}
Estas fórmulas producen dimensiones de frecuencia, energía y masa.
La velocidad y la longitud se obtienen mediante la respuesta de los
mismos canales, como precisan las identidades siguientes. Expresarlas
en segundos, julios o unidades de masa requiere las coordenatizaciones
correspondientes; esas coordenatizaciones no alteran los canales ni la sección
electrónica que se normaliza respecto al reloj.

\subsubsection{Identidad entre velocidad constitutiva y velocidad del reloj}

En el dominio \(x>|y|\), los dos canales satisfacen \(0<q_\pm<1\).
Sobre la realización de dos hojas, con proyectores de rango uno
\(P_+,P_-\), se define la respuesta temporal
\[
 \mathsf V=(I-T^{90})(I-T^{120})^{-1}
          =r_+P_++r_-P_-,
 \qquad r_\pm=\frac{1-q_\pm^{90}}{1-q_\pm^{120}}>0.
\]
El inverso existe porque ambos autovalores de \(T^{120}\) son menores
que uno. La lectura simétrica de la respuesta es
\[
 \mathcal E=
 \log[(1-q_+^{90})(1-q_-^{90})]
 -\log[(1-q_+^{120})(1-q_-^{120})]
 =\log(r_+r_-).
\]
Sean \(u_S,u_C,u_T\) las bases positivas de acción, velocidad y
tiempo, y \(u_L=u_Cu_T\) la base de longitud. Las lecturas
cinemática y constitutiva son
\[
 c_{\rm int}=e^{-\mathcal E}u_C,\qquad
 c_{\rm vac}=(r_+r_-)^{-1}u_C.
\]

\begin{proposition}[Realización común de propagación y masa absoluta]
\label{vi:dep:velocidad-comun}
En una misma coordenatización dimensional,
\[
 c_{\rm int}=c_{\rm vac}=\widehat c\,u_C,
 \qquad \widehat c=(\det\mathsf V)^{-1}=(r_+r_-)^{-1}.
\]
Para \(t_0=\tau_0u_T\), con \(\tau_0>0\), la longitud elemental es
\[
 \ell_0=c_{\rm int}t_0=\tau_0\widehat c\,u_L.
\]
La masa cronológica de \eqref{vi:dep:reloj-absoluto} conserva, por
tanto, la forma
\[
 m_{\rm clk}=
 \frac{2\pi\,\eta_{\rm ret}\,10^{-34}}
      {108\,\tau_0\,\widehat c^{\,2}}\,
 \frac{\mathcal U_S}{u_Tu_C^2}.
\]
\end{proposition}
\begin{proof}
La identidad \(\mathcal E=\log(r_+r_-)\) da
\(e^{-\mathcal E}=(r_+r_-)^{-1}\); el rango uno de cada hoja da
\(\det\mathsf V=r_+r_-\). El intercambio de hojas permuta los
factores y conserva su producto. La longitud se obtiene multiplicando
la misma velocidad por \(t_0\). Sustituir estas expresiones y
\eqref{vi:dep:hbar} en \eqref{vi:dep:reloj-absoluto} produce la
fórmula másica. El factor \(\mathcal U_S/(u_Tu_C^2)\) tiene
dimensión de masa; no interviene en la selección de los canales.
\end{proof}

Para una ruta homogénea de \(n>0\) transiciones, las lecturas aditivas
\(T(\gamma)=nt_0\) y \(L(\gamma)=n\ell_0\) dan
\(L(\gamma)/T(\gamma)=c_{\rm vac}\). Si el canal varía por
arista, se suman sus longitudes; el cociente es entonces la velocidad
media de esas aristas, no necesariamente una velocidad local común.

La igualdad de secciones permanece bajo cambios de coordenatización. Si
\(u_C'=du_C\), \(u_T'=\eta u_T\), con \(d,\eta>0\), entonces
\[
 \widehat c'=\widehat c/d,\qquad
 \tau_0'=\tau_0/\eta,\qquad u_L'=d\eta u_L,
 \qquad
 \tau_0'\widehat c'u_L'=\ell_0.
\]
En particular, elegir la base adaptada \(u_C'=c_{\rm vac}\)
hace \(\widehat c'=1\), pero la coordenada del lector en la
coordenatización interna sigue siendo \((r_+r_-)^{-1}\).

\subsection{Dos lectores de la trayectoria electrónica y evaluación beta}
\label{vi:dep:beta}

La trayectoria central empieza en \((r,c)=(5,5)\), con
orientación \((h,v)=(1,1)\). TRIT repite el ciclo \((+1,-1,0)\).
En el primer régimen se lee \(r+c\) y después se avanza
\(c\mapsto1+((c-1+h)\bmod9)\); en el segundo se lee \(rc\) y
después se avanza \(r\mapsto1+((r-1+v)\bmod9)\); en el régimen
cero se lee \(9\), sin desplazamiento. Se expresa la clase módulo
tres de la raíz digital de la lectura. Tras los pasos \(27,54,81\),
se invierten simultáneamente \(h,v\). La representación de \(108\)
pasos conserva la trayectoria completa y da
\begin{equation}
 a=100000220,\quad b=120200000,\quad
 \gamma_e=(a^3b^3)^2,
 \quad\tau_e=(+,+,+,-,-,-,+,+,+,-,-,-).
 \label{vi:dep:trayectoria-electron}
\end{equation}
La identidad se comprueba recorriendo los nueve instantes de cada
ventana; al tercer bloque la inversión cambia \(a\) por \(b\),
y el segundo recorrido de \(54\) pasos repite la representación.
El estado conserva asimismo sus coordenadas de memoria.

El primer lector mide desacuerdos cíclicos de la representación ternaria:
\[
 D_k=\#\{t\in\ZZ/108\ZZ:\gamma_e(t)\ne\gamma_e(t+k)\}.
\]
Como la palabra expresada tiene período \(54\), cada desacuerdo
aparece con su trasladado por \(54\); todos los \(D_k\) son pares.
En la mitad \(a^3b^3\), el recuento de desacuerdos para
\(k=1,3,4,27,36\) da respectivamente \(27,28,34,24,16\).
Al duplicarlo,
\[
 (D_1,D_3,D_4,D_{27},D_{36})=(54,56,68,48,32).
\]
Por tanto,
\begin{equation}
 K_D=\left(D_{27}+D_{36},D_1,\frac{D_4-D_3}{2}\right)
     =(80,54,6).
 \label{vi:dep:KD}
\end{equation}

El segundo lector actúa sobre la orientación de las doce ventanas.
Definamos, con índices cíclicos,
\[
 C_k=\sum_{r=0}^{11}\tau_r\tau_{r+k},\quad
 A_\ell=\sum_{r=0}^{11}\left(\sum_{j=0}^{\ell-1}\tau_{r+j}\right)^2,
 \quad P_6=\sum_{r=0}^{5}\tau_r\tau_{r+6}.
\]
La combinación entera primitiva orientada que cancela la parte diagonal
de \(A_3,A_4\) es \(\Omega=4A_3-3A_4\): sus coeficientes
satisfacen \(3u+4v=0\), \(\gcd(u,v)=1\), y se fija \(u>0\).
Expandir el cuadrado prueba
\[
 A_\ell=\ell C_0+
     2\sum_{k=1}^{\ell-1}(\ell-k)C_k,
 \qquad\Omega=-2C_1-4C_2-6C_3.
\]
Para \(\tau_e\), se obtiene
\(C_0=12,C_1=4,C_2=-4,C_3=-12\),
\(A_3=44,A_4=32,P_6=6\). De aquí,
\begin{equation}
 K_\Omega=(\Omega,54,P_6)=(80,54,6)=K_D.
 \label{vi:dep:KOmega}
\end{equation}
El segundo componente conserva el semirretorno de \(108\) pasos;
no se introduce por la cancelación diagonal. La igualdad de ambos
lectores se ha probado sobre la trayectoria central especificada;
no se extiende por definición a todas las colecciones del atlas.

\begin{definition}[Coordenada electrónica posterior al retorno]
La composición electrónica completa es
\begin{equation}
 \varepsilon_e^\beta
  =\varepsilon_e^{(0)}
      R_{\rm act}^{80-54\alpha+6\Delta_4}.
 \label{vi:dep:electron-beta}
\end{equation}
\end{definition}
La positividad de \eqref{vi:dep:formula-basal} y
\eqref{vi:dep:Ract} permite definir esta potencia mediante el
logaritmo real. El triple entero procede indistintamente de los dos
lectores centrales acabados de calcular. La evaluación relativa
preserva la constante reducida de Planck mediante la razón de sus
secciones, y la realización absoluta utiliza además
\eqref{vi:dep:reloj-absoluto} y la coordenatización dimensional declarada.

En las tablas heredadas, \(\mathfrak e_0\) denomina la
coordenada basal \(\varepsilon_e^{(0)}\) en la coordenatización allí
especificada. La sección \(\hbar_{\rm int}\) usada para una
evaluación de supervivencia debe declararse en la misma coordenatización:
para la realización posterior al retorno del presente desarrollo
es \(\hbar_{\rm ret}\). Las filas históricas conservan la
convención declarada en su tabla de origen y no se recalculan mediante un
cambio silencioso de esta sección.

Si \(\mathcal U_E\) es la base de energía de la evaluación
electrónica, \(E_e=\varepsilon_e^\beta\mathcal U_E\). Al emplear
el cuanto cronológico como base, el coeficiente es
\[
 \widehat\varepsilon_e
   =\varepsilon_e^\beta\frac{\mathcal U_E}{E_{\rm clk}},
 \qquad E_e=\widehat\varepsilon_e E_{\rm clk}.
\]
La igualdad conserva el objeto y transforma su coordenada. Por ello
el número que se presenta en una coordenatización de MeV no puede multiplicarse
sin conversión por cualquier otro cuanto de energía. El electrón
conserva su construcción estructural al expresarse respecto a una
escala cronológica: no se sustituye por la masa de ese cuanto.

La comparación metrológica se aplica al resultado de esta composición,
con la versión de \(\alpha\), la sección de acción, el lector
\(\Delta_4\) y la coordenatización de unidades declaradas. En particular,
\(d_4\) y \(\Delta_4\) no son intercambiables: reemplazarlos
simultáneamente en el factor basal y en el exponente cambia el
resultado por la razón positiva
\[
 \frac{1+15\Delta_4}{1+15d_4}
 R_{\rm act}^{6(\Delta_4-d_4)}>1.
\]
Las fórmulas conservadas permiten localizar esa diferencia sin
atribuirla a una masa externa ni seleccionar de nuevo la trayectoria.
